\endinput
%------------------------------------------------------------------------------%
\begin{question}
\end{question}

%------------------------------------------------------------------------------%
\begin{resolution}{Solução}
\end{resolution}

%------------------------------------------------------------------------------%
\endinput

\begin{frame}{Exemplo -- Escolhendo a melhor expansão}
Considere
\[
A=
\begin{pmatrix}
1&2&0&0\\
0&3&0&0\\
2&1&4&5\\
0&0&6&7
\end{pmatrix}.
\]
\medskip

A segunda coluna não possui muitos zeros, enquanto a segunda
linha possui três zeros. Portanto, é conveniente expandir pela
segunda linha:
\[
\det(A)
=
3C_{22}.
\]
\medskip

Como $2+2$ é par,
\[
C_{22}=M_{22}.
\]
\end{frame}

\begin{frame}{Exemplo -- Escolhendo a melhor expansão}
Temos
\[
M_{22}
=
\begin{vmatrix}
1&0&0\\
2&4&5\\
0&6&7
\end{vmatrix}.
\]
\medskip

Expandindo pela primeira linha:
\[
M_{22}
=
\begin{vmatrix}
4&5\\
6&7
\end{vmatrix}
=28-30=-2.
\]
\medskip

Portanto,
\[
\det(A)=3(-2)=-6.
\]
\medskip

Assim,
\[
\boxed{\det(A)=-6}.
\]
\end{frame}

